% Service properties of the CURRENT legacy_l10 Q planner and local ranking.
% Additive proof module; no executable or sealed empirical evidence is changed.
% Intended for the method chapter, before empirical evaluation.
\section{Tính chất của bước chọn Q và phục hồi dịch vụ}
\label{sec:formal-current-service}

Phần này tách ba câu hỏi: Q có thể di chuyển theo mô hình đường hay không;
nhu cầu riêng có làm đổi bản tin mạng hay không; và tập POI nhận về có đủ
để trả đúng top-$k$ hay không. Các tính chất này không thay thế định lý
Geo-I hoặc chứng minh phương pháp vượt đối chứng. Chúng áp dụng cho cấu
hình hiện tại: $K=5$, bảng chọn Q có $L_{\mathrm{plan}}=10$, phản hồi
có $L=30$ và đáp án local có $k=5$.

\subsection{Khả thi theo đường có hướng}
\label{sec:formal-service-reachability}

Giữ graph công khai $G=(V,E)$ cố định. Mỗi cạnh có độ dài không âm và tốc
độ dương; $T_G(v,w)$ là thời gian đường đi ngắn nhất có hướng, bằng
$+\infty$ khi không tới được. Cạnh nối hai trạng thái trùng tọa độ có thể
có chi phí 0. Miền $V_c\subseteq V$ là thành phần liên thông mạnh công
khai không rỗng được chọn trước chuyến. Gọi $v_{t,j}$ là \emph{trạng thái
làn--độ tiến nội bộ} của track $j$, còn $g(v_{t,j})$ là tọa độ Q gửi đi.
Với lịch công khai tăng nghiêm ngặt, miền chọn là
\begin{equation}
 \mathcal A_{t,j}=\begin{cases}
 V_c,&\text{sự kiện đầu},\\
 \{v\in V_c:T_G(v_{t-1,j},v)\leq\Delta t\},&\text{các sự kiện sau},
 \end{cases}
 \qquad \Delta t=t-t_{\mathrm{prev}}>0.
 \label{eq:formal-service-buckets}
\end{equation}

\noindent\textbf{Tính chất 1 (khả thi từng track).}
Nếu mọi bước chọn và thay trạng thái giữ $v_{t,j}\in\mathcal A_{t,j}$,
mỗi track có một đường đi trong $G$ giữa hai sự kiện liên tiếp với thời
gian không quá $\Delta t$.

\noindent\emph{Chứng minh.}
Miền đầu không rỗng. Ở bước sau, trạng thái trước vẫn thuộc $V_c$ và
$T_G(v_{t-1,j},v_{t-1,j})=0$, nên miền chọn không rỗng. Greedy chọn trong
miền; gộp trạng thái chỉ lấy tập con của miền; thay một track vẫn dùng
miền của track đó. Bước tiến tới mục tiêu chọn một trạng thái tới được
có cùng signature POI; bước slack cũng chỉ xét trạng thái tới được.
Quy nạp cho tính chất. Nối các đường giữa sự kiện và cho phép chờ bù thời
gian tạo một lộ trình theo mô hình graph. \hfill$\square$

Đây là tính chất của \textbf{trạng thái nội bộ}, không phải chứng minh
mọi phép suy track từ tọa độ đều khả thi. Hai làn khác nhau có thể trùng
tọa độ; phép truy hồi POI bằng trạng thái gần tọa độ nhất có thể chọn một
làn khác. Vì vậy chỉ khẳng định tồn tại cách gán trạng thái khả thi cho
chuỗi tọa độ đã sinh. Xáo thứ tự Q không xóa khả năng đối thủ nối track
bằng hình học. Các track cũng có thể trùng trạng thái; không có bảo đảm
tránh va chạm hoặc $K$ vị trí khác nhau.

Graph dùng tốc độ free-flow và danh mục trạng thái rời rạc công khai;
vòng hiện tại dùng khoảng chia không quá40 m và tốc độ không quá8 m/s.
Điều này không xác nhận đèn tín hiệu, tắc đường, gia tốc hoặc đổi làn
trong giao thông thật. Phép tính đường ngắn nhất lý tưởng cho tính chất
trên; kiểm tra SciPy/float chỉ xác nhận mô hình số đang chạy, chưa chứng
nhận sai số đối với đường liên tục.\footnote{Miền làn công khai:
\path{data/lane_states.py}; graph có hướng và phép tới được:
\path{evaluation/lane_travel.py}; các bước giữ miền:
\path{benchmark/engines/quotient_cover.py},
\path{benchmark/engines/progress_cover.py},
\path{benchmark/engines/slack_progress.py}.}

\subsection{Mục tiêu độ phủ và ý nghĩa chính xác của slack}
\label{sec:formal-service-cover}

Tại một sự kiện và lịch sử đã bảo vệ cố định, gọi $S_{10}(v)$ là hợp các
ID POI trong signature nearest của bộ chọn tại trạng thái $v$. Bảng này
theo phép ánh xạ tọa độ của dịch vụ, không được lấy trực tiếp từ GPS thật.
Trọng số $w_t(p)\geq0$ được tính từ belief trên lưới công khai. Trong cấu
hình hiện tại, mô hình gốc tạo trọng số từ \textbf{reference nearest
top5}; wrapper giữ các trọng số đó và dùng signature top10 để tính độ
phủ. Không đồng nhất hai độ sâu này với phản hồi $L=30$.

Với một lựa chọn trạng thái cho mỗi track, đặt
\begin{equation}
 F_t(Q)=\sum_p w_t(p)\,
 \mathbf1\{p\in\textstyle\bigcup_{j=1}^K S_{10}(v_{t,j})\}.
 \label{eq:formal-service-cover}
\end{equation}
Đây là điểm độ phủ nearest theo belief xấp xỉ. Nó không phải Recall thật
của bốn purpose, cũng không là kỳ vọng đã hiệu chuẩn cho một người dùng.
Trạng thái lưới không có reference đóng góp khối trọng số 0; không tạo
POI hoặc điểm số giả để thay reference rỗng.

\noindent\textbf{Tính chất 2 (độ phủ có lợi ích biên giảm dần).}
Xem mỗi phần tử chọn là cặp có nhãn $(j,v)$, với
$v\in\mathcal A_{t,j}$. Khi đó $F_t(\varnothing)=0$, $F_t$ không giảm
và là hàm submodular. Cụ thể, nếu $A\subseteq B$ thì
\[
 \Delta(e\mid A)=\sum_{p\in S_{10}(e)\setminus S_{10}(A)}w_t(p)
 \ \geq\
 \Delta(e\mid B)\geq0.
\]
Mở rộng tập chỉ thêm POI; một POI đã được phủ không tạo lợi ích lần hai.
Ràng buộc lấy không quá một phần tử mỗi nhãn tạo một matroid phân hoạch
(partition matroid). Vì mỗi miền không rỗng và $F_t$ không giảm, có thể hoàn
thành lựa chọn thành đúng $K$ phần tử, một phần tử cho mỗi track. Hai nhãn
vẫn có thể chọn cùng trạng thái đường.

\noindent\textbf{Hệ quả (cận greedy kế thừa).}
Với oracle lợi ích biên chính xác trên toàn bộ các miền không rỗng,
greedy toàn cục đang dùng đạt $F_t(Q_{\mathrm g})\geq F_t(Q^*)/2$, trong
đó $Q^*$ tối ưu \emph{cho cùng sự kiện, belief và miền chọn}.
Đây là cận greedy chuẩn dưới ràng buộc matroid~\cite{calinescu2011matroid}, không là đóng góp
nguyên lý mới. Có thể chứng minh trực tiếp như sau. Gọi $a_i$ là phần tử
greedy chọn ở bước $i$, $S_i=\{a_1,\ldots,a_i\}$ và $o_i$ là phần tử
của $Q^*$ thuộc cùng track với $a_i$. Khi chọn $a_i$, track đó chưa được
chọn, nên $o_i$ còn khả thi. Greedy cho
$\Delta(o_i\mid S_{i-1})\leq\Delta(a_i\mid S_{i-1})$. Do tính không
giảm và submodular,
\begin{equation}
 \begin{split}
 F_t(Q^*)&\leq F_t(S_K\cup Q^*)\\
 &\leq F_t(S_K)+\sum_{i=1}^K\Delta(o_i\mid S_K)\\
 &\leq F_t(S_K)+\sum_{i=1}^K\Delta(o_i\mid S_{i-1})
 \leq2F_t(S_K).
 \end{split}
 \label{eq:formal-service-greedy-half}
\end{equation}

Gộp những \emph{signature có cùng thứ tự} giữ độ phủ và đại diện có
tie-break công khai tốt nhất tại bước hiện tại. Thay một track được nhận
chỉ khi toàn mục tiêu tăng; tiến tới mục tiêu giữ signature nên giữ
$F_t$. Giới hạn tối đa3 lần thay không chứng minh đã tới tối ưu cục bộ.
Gọi $Q_{\mathrm{pre}}$ là lựa chọn trước slack và
$\delta_F=\texttt{utility\_slack}$. Mỗi thay đổi slack được kiểm tra lại
với cùng floor $F_t(Q_{\mathrm{pre}})-\delta_F$, do đó trong số học lý
tưởng
\begin{equation}
 F_t(Q_{\mathrm{final}})\geq F_t(Q_{\mathrm{pre}})-\delta_F
 \geq\tfrac12F_t(Q^*)-\delta_F.
 \label{eq:formal-service-slack-floor}
\end{equation}
Ở bước đầu hoặc khi slack bằng0, bước slack không đổi lựa chọn. Giá trị
hiện tại $\delta_F=0{,}03$ là mất tối đa0,03 \emph{điểm surrogate cùng
bước}, không phải3\% sai số Recall hay xác suất lộ vị trí. Mã chấp nhận
floor với tolerance $10^{-12}$; đó là kiểm tra giá trị float ghi nhận,
không là chứng chỉ sai số $10^{-12}$ so với số thực. Cận $1/2$ cần oracle
chính xác, không tự áp như chứng chỉ executable. Trạng thái sau slack
làm đổi miền bước sau, nên không suy cận độ phủ cho toàn chuyến hoặc
chất lượng người dùng.\footnote{Trọng số reference5 được giữ trong
\path{benchmark/response_aware_belief.py}; nguồn resource:
\path{experiments/future_sumo_eval.py}, hàm \texttt{native\_resources}.
Greedy, độ phủ và full-action guard:
\path{benchmark/engines/service_cover.py},
\path{benchmark/engines/fair_cover.py}; factory active:
\path{benchmark/engines/public_service_planner.py}.}

\subsection{Không thêm kênh request theo nhu cầu riêng}
\label{sec:formal-service-noninterference}

Gọi $\psi$ là chuỗi nhu cầu local: purpose, category, bán kính, đích và
top-$k$. Giữ cùng trace vị trí $X$, ngữ cảnh công khai $\mathcal C$, lịch
mở/đóng, admission, sự kiện, mọi loại được request và $L$ cố định. Giữ
cùng chiến lược máy chủ. Giả sử $\psi$ chỉ đi vào local ranking, không
đổi Q, lịch đọc/gửi, expiry, reset hoặc retry. Với
$\mathsf W(X,\psi;\mathcal C)$ là transcript \emph{mạng logic}, ta có
\begin{equation}
 \mathsf W(X,\psi;\mathcal C)
 \overset{\mathrm d}{=}
 \mathsf W(X,\psi';\mathcal C)
 \quad\text{cho mọi }\psi,\psi'.
 \label{eq:current-purpose-noninterference}
\end{equation}

\noindent\emph{Chứng minh.}
Ghép cùng randomness của REM, planner, thứ tự Q và máy chủ. Ở sự kiện
đầu, các đầu vào tạo bản tin giống nhau. Nếu các prefix mạng giống nhau,
lịch sử dùng cho sự kiện tiếp theo giống nhau; nhu cầu local không sửa
chúng. Mỗi Q request mọi category công khai với cùng $L$, nên cả request
và response tiếp theo giống nhau. Quy nạp cho bình đẳng từng bản tin
trên cùng randomness, mạnh hơn bình đẳng phân bố. \hfill$\square$

API mạng không nhận \texttt{QuerySpec}; hàm trả lời local không gọi GPS
supplier của Geo-I, không fetch và không tiến đồng hồ cache. Điều kiện
này bảo vệ \emph{nội dung request có điều kiện theo cùng vị trí/lịch}.
Nó không nói intent độc lập với tuyến tới bệnh viện, giờ bật ứng dụng,
account/IP, click hoặc đặt chỗ. Logical-clock proof cũng không chứng
minh latency thực độc lập nhu cầu nếu CPU/cache contention hay lỗi local
ảnh hưởng packet. Nếu thiếu đáp án làm ứng dụng gửi request mới theo
nhu cầu thì điều kiện đã bị phá.\footnote{Ranh giới API:
\path{benchmark/query_purpose.py},
\path{benchmark/budgeted_geoi_lbs.py}; kiểm tra local scope chung cho
hai cách dùng cache: \path{benchmark/versioned_static_geoi_lbs.py}.}

\subsection{Khi nào local top-$k$ là đáp án chính xác?}
\label{sec:formal-service-local-exactness}

Giữ vị trí local đúng trong mô hình: $v$ là trạng thái gần GPS hiện tại
theo phép ánh xạ công khai. Giữ một catalogue/version và nhu cầu $\psi$.
Với category đã chọn, các điểm số hiện thực là $D_G(v,p)$ cho nearest,
$T_G(v,p)$ cho fastest, $D_G(v,p)$ với điều kiện $D_G(v,p)\leq r$ cho
within-radius, và
$D_G(v,p)+D_G(p,a)-D_G(v,a)$ cho detour. Chỉ nhận điểm số hữu hạn.
Nếu đích $a$ không tới được thì detour có reference rỗng. Bán kính là
khoảng cách đường có hướng, còn fastest dùng tốc độ free-flow, không
phải thời gian giao thông đo thực. Tie được phá bằng ID POI công khai.

Gọi $E_\psi(v)$ là miền POI hợp điều kiện và $\prec$ là thứ tự toàn phần
theo điểm số rồi ID. Đặt
\begin{equation}
 R=\operatorname{Top}_k(E_\psi(v)),\quad
 r_R=|R|=\min(k,|E_\psi(v)|),\quad
 P(A)=\operatorname{Top}_k(A\cap E_\psi(v)),
 \label{eq:formal-service-reference}
\end{equation}
trong đó $A$ là tập ID nhận hợp lệ. Với dịch vụ tĩnh hiện tại, tính hợp lệ
yêu cầu đúng catalogue/version. Nếu xét availability động, phải bổ sung
trạng thái \emph{hiện tại đã biết}; chỉ có metadata không đủ xác nhận POI
đang hoạt động. Detour cần đích local đã biết; dùng endpoint thật trong
evaluator không chứng minh thiết bị suy được đích từ quá khứ.

\noindent\textbf{Tính chất 3 (thu hồi reference).}
Khi $r_R>0$, nếu đáp án và reference dùng \emph{cùng thứ tự và miền đúng},
\begin{equation}
 R\cap P(A)=R\cap A,\qquad
 \operatorname{Recall}@k(A)=\frac{|R\cap A|}{r_R}.
 \label{eq:formal-service-recall-identity}
\end{equation}
Một $p\in R\cap A$ có hạng toàn miền không quá $k$; bỏ POI ngoài $A$
không thể làm hạng đó lớn hơn. Do đó $p$ vẫn nằm trong $P(A)$. Chiều
ngược hiển nhiên vì $P(A)\subseteq A$. Đặc biệt, $R\subseteq A$ là điều
kiện cần và đủ cho $P(A)=R$ và Recall bằng1. Không cần tải toàn catalogue,
nhưng không có bảo đảm bộ chọn Q luôn thu hồi đủ $R$.

Reference có thể ít hơn5 POI. Khi $r_R=0$, Recall là N/A và số reference
rỗng phải được báo riêng. Khi reference có nghĩa nhưng câu trả lời rỗng,
Recall bằng0; không đổi mẫu số thành kích thước reference của vị trí
ước lượng. Completion đủ số lượng cũng không đồng nghĩa Recall đúng:
có thể trả5 POI hợp điều kiện nhưng cả5 đều ngoài reference top5.
\footnote{Điểm số và tie-break:
\path{benchmark/query_purpose.py}, lớp \texttt{MultiPurposeRoadRanking};
reference/mẫu số của vòng hiện tại:
\path{experiments/qplanner_study_20261006_v2.py},
lớp \texttt{UtilityEvaluator}.}

\subsection{Tăng L: điều kiện đơn điệu và phản ví dụ}
\label{sec:formal-service-depth-monotonicity}

Giữ \emph{đúng cùng Q}, graph, phép truy cập tọa độ, catalogue tĩnh và
thứ tự dịch vụ. Với mỗi Q/category, top20 là prefix của top30, nên hợp
ID sau khử trùng thỏa $A_{20}\subseteq A_{30}$; có thể bằng nhau khi ít
POI. Nếu giữ cùng reference và local ranking đúng như Tính chất 3,
\begin{equation}
 \operatorname{Recall}@k(A_{30})-
 \operatorname{Recall}@k(A_{20})
 =\frac{|R\cap(A_{30}\setminus A_{20})|}{r_R}\geq0.
 \label{eq:formal-service-depth-monotonicity}
\end{equation}
Do đó các trung bình với cùng tập reference có nghĩa và cùng trọng số
không âm cũng không giảm. Không suy tăng Recall nghiêm ngặt, tăng
privacy hoặc ưu thế tại cùng bandwidth; tăng $L$ có giá phản hồi. Thay
$L_{\mathrm{plan}}$ làm Q khác đi thì lập luận giữ-cùng-Q không còn áp dụng.

\textbf{Sai vị trí local có thể phá tính đơn điệu, dù sai giống nhau ở
hai độ sâu.} Trên đường hai chiều, vị trí thật ở0; năm POI $a_1,\ldots,a_5$
ở1,2,3,4,5, còn $b$ ở10. Với $k=5$, reference thật là năm $a_i$.
Nếu pool nhỏ chứa chúng, local top5 trả cả năm dù ước lượng vị trí ở10.
Khi pool lớn thêm $b$, local ranking sai đưa $b$ lên đầu, đẩy một $a_i$
ra: Recall giảm từ1 xuống$4/5$. Đây là phản ví dụ về tập ứng viên lồng
nhau. Để có đúng prefix20/30, trên một graph có hướng có thể thêm15 POI
xa cả vị trí thật và vị trí ước lượng, nhưng nằm sau năm $a_i$ và trước
$b$ trong thứ tự từ Q. Chẳng hạn, ba trạng thái nguồn có các cạnh trực
tiếp tới POI với chi phí cho đúng ba thứ tự đó; các cạnh quay lại đủ dài
giữ những thứ tự và làm graph liên thông mạnh. Pool20 có năm $a_i$ và15
POI bổ sung; pool30 thêm $b$. Thu hồi thêm reference vẫn không giảm,
nhưng Recall \emph{sau xếp hạng sai} có thể giảm.

\textbf{Cache không tự phá tính đơn điệu.} Nếu hai cache bắt đầu giống
nhau, cùng lịch reset/expiry/version và cùng quy tắc hiệu lực, hợp các
phản hồi prefix giữ $A_{20}(t)\subseteq A_{30}(t)$ bằng quy nạp. Với
availability, trạng thái của ID chung còn phải nhất quán ở cùng epoch.
Ngược lại, so cache tích lũy L20 giữ POI $a$ từ quá khứ với L30 current-only
không còn $a$ thì các pool không lồng nhau. Eviction khác nhau cũng có
thể phá điều kiện này. Trạng thái cũ bị hết hạn trở thành unknown; không
được xem unknown là đang hoạt động hoặc là chắc chắn không hoạt động.

\textbf{Khác đồng hồ hoặc phiên bản không là so prefix.} Nếu phản hồi
L20 lấy trước cập nhật catalogue còn L30 lấy sau khi POI $a$ bị bỏ,
$a$ có thể nằm ở pool20 mà không có ở pool30. Reference hiện tại cũng
có thể đổi khi người dùng di chuyển hoặc availability đổi. Vì thế phải
giữ cùng thời điểm, miền, catalogue và chính sách cache khi kiểm tra
độ sâu; không lấy số liệu hai trạng thái dịch vụ khác nhau để suy một
bất đẳng thức từng event.

Các chứng minh trên giải thích điều kiện thiết kế và cách đọc phép đo.
Chúng không chứng minh posterior chính xác, GPS local không sai, dịch vụ
động có đủ trạng thái hiện hành hoặc mọi scenario đã được bảo vệ. Những
điểm đó được đánh giá riêng bằng benchmark và các diagnostic có phạm vi
nguồn, mẫu số và chi phí được ghi rõ.
